% Part: sets-functions-relations
% Chapter: size-of-sets
% Section: pairing-alt
% Hindi translation (hi-Deva-IN) of Open Logic commit 9620cc73f9c8e0ad003c514a5d3748f29611c4c0.

\documentclass[../../../include/open-logic-section]{subfiles}

\begin{document}

\olfileid[hi]{sfr}{siz}{pai-alt}

\olsection{एक वैकल्पिक युग्मन फलन}

\begin{explain}
$\Nat^2$ की कुछ अन्य गणनाएँ ऐसी हैं जिनके प्रतिलोम ज्ञात करना अधिक सरल है।
उनमें से एक यह है। गणना को सरणी में देखने के स्थान पर, आरंभ में रिक्त स्थानों
से संबद्ध धन पूर्णांकों की सूची से आरंभ कीजिए। कल्पना कीजिए कि हम इन स्थानों
को क्रमशः $\tuple{n,m}$ युग्मों से इस प्रकार भरते हैं। पहले उन युग्मों को लें
जिनके पहले स्थान पर~$0$ है (अर्थात $\tuple{0,m}$ रूप के युग्म)। पहले युग्म
(अर्थात $\tuple{0,0}$) को पहले रिक्त स्थान में रखिए, फिर एक रिक्त स्थान छोड़िए,
दूसरे युग्म (अर्थात $\tuple{0,1}$) को अगले रिक्त स्थान में रखिए, फिर एक स्थान
छोड़िए, और इसी प्रकार आगे बढ़ते जाइए। अब हमारी गणना का (अधूरा) आरंभ ऐसा
दिखता है:
\[\small
\begin{array}{@{}c c c c c c c c c c c@{}}
\mathbf 1 & \mathbf 2 & \mathbf 3 & \mathbf 4 & \mathbf 5 & \mathbf 6 & \mathbf 7 & \mathbf 8 & \mathbf 9 & \mathbf{10} & \dots \\ \\
\tuple{0,0} &  & \tuple{0,1} &  & \tuple{0,2} &  & \tuple{0,3} & & \tuple{0,4} &  & \dots \\
\end{array}
\]
अब भी रिक्त बचे स्थानों में $\tuple{1,m}$ युग्मों के साथ यही प्रक्रिया
दोहराइए और फिर प्रत्येक दूसरे रिक्त स्थान को छोड़ते जाइए:
\[\small
\begin{array}{@{}c c c c c c c c c c c@{}}
\mathbf 1 & \mathbf 2 & \mathbf 3 & \mathbf 4 & \mathbf 5 & \mathbf 6 & \mathbf 7 & \mathbf 8 & \mathbf 9 & \mathbf{10} & \dots \\ \\
\tuple{0,0} & \tuple{1,0} & \tuple{0,1} &  & \tuple{0,2} & \tuple{1,1} & 
\tuple{0,3} & & \tuple{0,4} &  \tuple{1,2} & \dots \\
\end{array}
\]
इसी प्रकार $\tuple{2,m}$, $\tuple{3,m}$ इत्यादि युग्म भरिए। इस प्रकार हमारी
पूर्ण गणना का आरंभ ऐसा होता है:
\[\small
\begin{array}{@{}cc c c c c c c c c c@{}}
\mathbf 1 & \mathbf 2 & \mathbf 3 & \mathbf 4 & \mathbf 5 & \mathbf 6 & \mathbf 7 & \mathbf 8 & \mathbf 9 & \mathbf{10} & \dots \\ \\
\tuple{0,0} & \tuple{1,0} & \tuple{0,1} & \tuple{2,0}  & \tuple{0,2} & 
\tuple{1,1} & \tuple{0,3} & \tuple{3,0}  & \tuple{0,4} &  \tuple{1,2} & \dots \\
\end{array}
\]
यदि हम ऊपर की सरणी के खानों को इस गणना के अनुसार क्रमांक दें, तो हमें कोई
सुव्यवस्थित आड़ी-तिरछी रेखा नहीं, बल्कि यह विन्यास मिलेगा:
\[
\begin{array}{ c | c | c | c | c | c | c | c }
& \mathbf 0 & \mathbf 1 & \mathbf 2 & \mathbf 3 & \mathbf 4 & \mathbf 5 & \dots \\
\hline
\mathbf 0 & 1 & 3 & 5 & 7 & 9 & 11 & \dots \\
\hline
\mathbf 1 & 2 & 6 & 10 & 14 & 18 & \dots & \dots \\
\hline
\mathbf 2 & 4 & 12 & 20 & 28 & \dots & \dots & \dots \\
\hline
\mathbf 3 & 8 & 24 & 40 & \dots & \dots & \dots & \dots \\
\hline
\mathbf 4 & 16 & 48 & \dots & \dots & \dots & \dots & \dots \\
\hline
\mathbf 5 & 32 & \dots & \dots & \dots & \dots & \dots & \dots \\
\hline
\vdots & \vdots & \vdots & \vdots & \vdots & \vdots & \vdots & \ddots\\
\end{array}
\]
हम देख सकते हैं कि पंक्ति~$0$ के युग्म हमारी गणना में विषम क्रमांक वाले
स्थानों पर हैं, अर्थात युग्म $\tuple{0,m}$ स्थान $2m+1$ पर है। दूसरी पंक्ति
के युग्म $\tuple{1,m}$ ऐसे स्थानों पर हैं जिनके क्रमांक किसी विषम संख्या के
दुगुने हैं, विशेष रूप से $2 \cdot (2m+1)$। तीसरी पंक्ति के युग्म
$\tuple{2,m}$ ऐसे स्थानों पर हैं जिनके क्रमांक किसी विषम संख्या के चार गुने
हैं, अर्थात $4 \cdot (2m+1)$; और इसी प्रकार आगे। प्रत्येक पंक्ति में
$(2m+1)$ के गुणक $1$, $2$, $4$, $8$, \dots, ठीक~$2$ की घातें हैं:
$1= 2^0$, $2 = 2^1$, $4 = 2^2$, $8 = 2^3$, \dots\@ वस्तुतः, संबंधित
घातांक सदा विचाराधीन युग्म का पहला अवयव होता है। अतः युग्म $\tuple{n,m}$ के
लिए गुणक $2^n$ है। इससे हमें सामान्य सूत्र $2^n \cdot (2m+1)$ मिलता है।
लेकिन यह युग्मों का \emph{धन} पूर्णांकों में प्रतिचित्रण है, अर्थात
$\tuple{0,0}$ का स्थान~$1$ है। यदि हम स्थान~$0$ से आरंभ करना चाहें, तो हमें
परिणाम में से~$1$ घटाना होगा। इससे हमें निम्नलिखित मिलता है:
\end{explain}

\begin{ex}
निम्नलिखित रूप से दिया गया फलन $h\colon \Nat^2 \to \Nat$,
\[
h(n,m) = 2^n (2m+1) - 1
\]
प्राकृत संख्याओं के युग्मों के समुच्चय~$\Nat^2$ का युग्मन फलन है।
\end{ex}

\begin{explain}
तदनुसार, $\Nat^2$ की हमारी दूसरी गणना में युग्म $\tuple{0,0}$ का कूट
$h(0,0) = 2^0(2\cdot 0+1) - 1 = 0$ है; $\tuple{1,2}$ का कूट
$2^{1} \cdot (2 \cdot 2 + 1) - 1 = 2
\cdot 5 - 1 = 9$ है; और $\tuple{2,6}$ का कूट $2^{2} \cdot (2
\cdot 6 + 1) - 1 = 51$ है।
\end{explain}

कभी-कभी प्राकृत संख्याओं के युग्मों~$\Nat^2$ को इस शर्त के बिना कूटित करना
ही पर्याप्त होता है कि कूटन आच्छादक हो। ऐसे कूटनों के प्रतिलोम केवल आंशिक
फलन होते हैं।

\begin{ex}
निम्नलिखित रूप से दिया गया फलन $j\colon \Nat^2 \to \Nat^+$,
\[
j(n,m) = 2^n3^m
\]
एक !!a{injective} फलन $\Nat^2 \to \Nat$ है।
\end{ex}

\end{document}
